% Section 1: Introduction and programme.

\subsection{Two native formalisms}

Arithmetic expression geometry begins with a sequential arithmetic expression and
keeps, rather than discards, its evaluation history.  Paper~I establishes the
intrinsic classification of sequential trees, the marked spinal history, the
chirality word, and the projective semantics those histories induce
\cite{Yuan2026AEGFoundations}.  Paper~0 develops the geometry of the resulting
spaces, and Paper~V develops the six-state carrier on which three readings of one
construction are compared \cite{Yuan2026AEGPaperZero,Yuan2026AEGPaperFive}.

Process Geometry begins one level lower.  Its native object is an \emph{open finite
program}: a typed expression DAG with ordered holes, explicit source and occurrence
identity, explicit copy and discard, and retained construction history
\cite{AdvaProcessCore2026,ProcessGeometryCore2026}.  Composition grafts an actual argument program into
a hole; evaluation is one observation of the resulting process, not its identity.
One process carries three coupled readings --- construction $K$, causal development
$C$, and frontier $F$ --- and the theory's recurring statement is negative: the same
value, the same differential, or the same spectrum does not imply the same program.

The two formalisms meet at a single observation.  The iota carrier of Paper~V has
\emph{arity one}: it acts at one slot, and acting at one slot induces an order on
slots.  Once that order is available, the spine of a sequential arithmetic
expression and the left and right combs that generate it are not primitive data but
the single-spine case of a process order.  This paper states the general theory that
results and the equivalence with arithmetic expression geometry to which it leads.

\subsection{What is established, and what this paper adds}

Three bodies of work are imported and cited rather than reproved.

\begin{enumerate}
\item \emph{The carrier.}  Paper~V defines the six-state carrier
  $W=\Lambda^{1}E\oplus\Lambda^{2}E$ with its three aspect--opposite-relation
  planes, the complex unit $\Omega=TSR$, and the chirality sign law
  \cite{Yuan2026AEGPaperFive}.  The agreement of this carrier with the Adva iota
  frame is recorded in the carousel--iota overview note (note~40) as a
  correspondence; it is not a claim of Paper~V's appendix.  This paper uses the
  carrier as given and cites the correspondence where it uses it.
\item \emph{The combinatorics.}  Pushkarev--Byzov give the carousel of a plane
  binary tree, its five corner classes and the four-cycle that permutes them, and
  the order-two shadow on the corner counts \cite{PushkarevByzov2014}.  Donaghey's
  transform between general and binary Catalan trees supplies the right/left
  bijections and the factorisation $\sigma,\theta,\varphi$
  \cite{Donaghey1980,PushkarevByzov2014}.
\item \emph{The process corpus.}  Process Geometry supplies the native program,
  the three readings, cuts and slices, and their exact laws at declared revisions
  \cite{AdvaProcessCore2026,ProcessGeometryCore2026}.
\end{enumerate}

On that basis the paper states:

\begin{itemize}
\item a definition of the program process and of an evaluation schedule as a linear
  extension of its construction order (\S\ref{sec:processes});
\item the arity-one principle, which recovers the spine and the comb orders from a
  slot order (\S\ref{sec:processes});
\item the evaluation-order theorem: for every plane binary tree, the left-first
  postorder is the unique linear extension of the dependency order minimizing
  inversions against the corner order (\S\ref{sec:orders});
\item the carousel correspondence: up to the non-canonical alignment of
  Proposition~\ref{prop:screened-alignment}, the class-level four-cycle aligns with
  the four-cycle of the iota complex unit on the signed states of one plane; the
  count shadow is an involution; and
  both are realized on the planar screw module
  $V_{\mathrm{class}}\otimes V_{\mathrm{index}}$ by $C^{4}=S$
  (\S\ref{sec:carousel});
\item the arithmetic--process dictionary and the equivalence statement it supports
  (\S\ref{sec:equivalence}), with the boundaries collected in
  \S\ref{sec:boundaries}.
\end{itemize}

\subsection{Status discipline}

Every statement carries one of the repository's status labels.  In this paper the
labels are used as follows, and the summary table of
\S\ref{sec:boundaries} records the assignment statement by statement.

\begin{center}
\small
\begin{tabular}{@{}ll@{}}
\toprule
Label & Used for \\
\midrule
\texttt{PROVED} & a statement proved here or in a cited proof record \\
\texttt{IMPORTED TECHNICAL} & a companion technical report at a declared revision \\
\texttt{IMPORTED CLASSICAL} & an external theorem, with attribution \\
\texttt{DEFINITION} & an object or construction fixed by the paper \\
\texttt{CORRESPONDENCE} & a stated alignment with cited evidence, not a theorem \\
\texttt{COMPUTATIONALLY VERIFIED EXAMPLE} & a finite instance checked by a stored verifier \\
\texttt{STRUCTURAL PROPOSAL} & a reading or construction, not a theorem \\
\texttt{CONJECTURE} & a precise statement expecting proof or refutation \\
\texttt{OPEN PROBLEM} & a question with no settled answer \\
\bottomrule
\end{tabular}
\end{center}

Three warnings apply throughout.  First, a finite calibration is not a theorem.
Second, an alignment that requires a choice is stated as non-canonical.  Third, the
equivalence of \S\ref{sec:equivalence} is stated at the strength the dictionary
supports; the directions it does not support are listed as open problems rather than
as consequences.

\subsection{Relation to the earlier papers}

The dependency is one-way and is fixed by
\path{governance/00e-paper-VI-process-geometry-scope.md}:

\begin{center}
\begin{tabular}{@{}ll@{}}
\toprule
Paper & What this paper imports \\
\midrule
Paper I & sequential histories, spine, chirality word, mirror and reversal \\
Paper 0 & the expression-space language and the charge vocabulary \\
Paper V & the six-state carrier, $\Omega$, and the chirality sign law \\
\bottomrule
\end{tabular}
\end{center}

The paper proves no new statement inside the scope of those papers, and no proved
node of Papers~0--V may depend on it.  What it adds is the general process theory
and the dictionary; what it does not do is replace the sequential theory.  The
sequential case is recovered as the single-spine leaf of the process theory, and the
recovery is stated as a proposition rather than as a redefinition.

\subsection{Roadmap}

Section~\ref{sec:processes} fixes the native objects: the open program, the three
readings, the schedule, the arity-one principle, and the spine and comb orders.
Section~\ref{sec:orders} states and proves the canonical evaluation order and records
its finite calibration.  Section~\ref{sec:carousel} develops the carousel
correspondence and the screw carrier, including the refutation of the strong form.
Section~\ref{sec:equivalence} gives the typed dictionary and the equivalence
statements.  Section~\ref{sec:boundaries} collects the claim ledger, the open
problems, and the directions that would extend the paper.
